\section{Related Work}
\label{sec:related}

\paragraph{Token-concatenation VLMs.}
LLaVA-style models connect a pretrained vision encoder to an LLM through a
light projector and train on instruction data~\citep{liu2023llava,
liu2024improved}. Scaling this recipe mostly means scaling the visual token
count, which is precisely what makes inference expensive at high resolution.
Token-pruning and spatial-pooling methods such as FastV~\citep{chen2024fastv}
and locality-aware projectors~\citep{cha2024honeybee} compress the visual sequence
to a smaller token budget, but still keep visual tokens inside every backbone layer.
Our baselines cover both full-token concatenation and $32$-token pruning/pooling
under our exact data and compute budget so that comparisons are apples-to-apples.

\paragraph{Parameter-efficient and encoder-free adaptation.}
Encoder-free VLMs internalise visual processing into the LLM by duplicating
modality-specific weights~\citep{diao2024eve} or inserting low-rank
adapters~\citep{hu2021lora,vora2025} into the early layers and distilling a
vision tower. While this removes the external ViT, unfreezing copied layers
requires ratusan of millions of parameters, and permanently merging additive
adapters into a text model destroys arithmetic reasoning, which motivates a
multiplicative, context-gated form of adaptation.

\paragraph{Rotation-based injection.}
ROSE~\citep{rose2025} converts visual features into block-diagonal orthogonal
matrices and rotates linear inputs, so the backbone never sees a visual token.
The construction guarantees that activations preserve their norm, and it is
attractive because the text model remains intact. Its limitations---a static,
globally pooled rotation, dependence on an external tower, and exposure bias
under greedy decoding---are the starting point of this work. We keep the
orthogonal construction, make it dynamic and spatially addressable, remove the
tower, and add on-policy distillation.

\paragraph{Distillation for VLMs.}
Knowledge distillation transfers a teacher's output distribution or feature
geometry to a student~\citep{hinton2015distilling,park2019rkd}; sequence-level
and on-policy variants~\citep{kim2016sequence,gu2024minillm,agarwal2024gkd}
train on student-generated trajectories. Off-policy distillation on
ground-truth prefixes does not cover the states a model visits when it decodes
on its own, whereas mode-seeking reverse KL~\citep{gu2024minillm} or hard top-$1$
sequence targets~\citep{kim2016sequence} can destabilise termination under
multiplicative rotation. Token-level forward KL on student rollouts closes this
gap and is the second step of our framework.

\paragraph{Orthogonal parameterisation.}
The Cayley transform $R=(I-A)^{-1}(I+A)$ maps a skew-symmetric matrix to a
special orthogonal matrix~\citep{golub2013matrix} and has been used to
stabilise recurrent and convolutional models. We use it to parameterise
block-diagonal rotations that are cheap to generate, exactly orthogonal, and
composable with a learned diagonal rescaling.

\paragraph{Benchmarks.}
We evaluate on general perception (MMStar~\citep{chen2024mmstar}), visual
hallucination (HallusionBench~\citep{guan2024hallusionbench}), chart reasoning
(ChartQA~\citep{masry2022chartqa}), scene-text understanding
(OCRBench~\citep{liu2023ocrbench}), mathematical reasoning
(MathVista~\citep{lu2024mathvista}), and plane geometry
(Geometry3K~\citep{lu2021geometry3k}). Text retention is measured with GSM8K
\citep{cobbe2021gsm8k} and ARC-Challenge~\citep{clark2018arc}.
